\subsection{Green 函数}\label{sec:greens-functions}

在前两节中, 我们考察了常系数偏微分方程的基本解. 这样的基本解使我们能够求解在整个空间上提出的形如 $L(D)u=f$ 的问题. 但在实际应用中, 通常并不希望求解整个空间上的问题; 而是希望求解某个区域上的偏微分方程, 并在边界上施加一定条件. Green 函数是这种情形下基本解的对应物. 只有在很特殊的情形下才能显式求出它们. 尽管如此, Green 函数这一概念对于理论研究仍然有用. 目前尚无可用的方法来讨论偏微分方程 Green 函数的存在性, 唯一性和正则性, 因而本节讨论在很大程度上将是形式的.

\subsubsection{边值问题及其伴随}\label{subsec:boundary-value-problems-adjoints}

\begin{Definition}
\label{def:formal-adjoint-differential-operator}
设
\begin{equation}\label{eq:variable-coefficient-differential-operator}
L(\bm{x},D)=\sum_{|\alpha|\leq k}a_\alpha(\bm{x})D^\alpha
\end{equation}
为定义在 $\Omega$ 上的微分算子. 那么 $L$ 的\textbf{形式伴随}是由下式给出的算子
\begin{equation}\label{eq:formal-adjoint-differential-operator}
L^*(\bm{x},D)u=\sum_{|\alpha|\leq k}(-1)^{|\alpha|}D^\alpha(a_\alpha(\bm{x})u(\bm{x})).
\end{equation}

\end{Definition}

这一定义的重要性在于, 有
\begin{equation}\label{eq:formal-adjoint-test-function-identity}
(\phi,L(\bm{x},D)\psi)=(L^*(\bm{x},D)\phi,\psi)
\end{equation}
对每个 $\phi,\psi\in\mathcal{D}(\Omega)$ 成立. 若去掉紧支集假设, 则 \eqref{eq:formal-adjoint-test-function-identity} 不再成立; 分部积分会产生含有边界 $\partial\Omega$ 上积分的附加项. 不过, 若 $\phi$ 和 $\psi$ 在边界上满足一定限制, 这些边界项就会消失. 我们关注的情形是 $L$ 的阶数为偶数, $k=2p$, 并给定了关于 $\psi$ 的 $p$ 个线性齐次边界条件. 此时, 自然要寻找施加于 $\phi$ 的 $p$ 个边界条件, 使 \eqref{eq:formal-adjoint-test-function-identity} 成立. 这就引出了伴随边值问题的概念.

为使这一想法具体化, 先考察常微分方程的情形. 设
\begin{equation}\label{eq:even-order-ordinary-differential-operator}
L(x,D)u=\sum_{i=0}^{2p}a_i(x)\dfrac{\mathrm{d}^iu}{\mathrm{d}x^i}(x),\quad x\in(a,b).
\end{equation}
假设 $a_i\in C^i([a,b])$; 这保证 $L^*$ 的系数连续. 此外, 假设对 $x\in[a,b]$ 有 $a_{2p}(x)\ne0$. 对任意函数 $\phi,\psi\in C^{2p}[a,b]$, 计算得
\begin{equation}\label{eq:ode-formal-adjoint-boundary-expansion}
\begin{aligned}
&(\phi,L(x,D)\psi)-(L^*(x,D)\phi,\psi)\\
&\quad=\sum_{i=1}^{2p}\sum_{k=0}^{i-1}(-1)^kD^{i-k-1}\psi(b)D^k(a_i\phi)(b)\\
&\qquad-\sum_{i=1}^{2p}\sum_{k=0}^{i-1}(-1)^kD^{i-k-1}\psi(a)D^k(a_i\phi)(a).
\end{aligned}
\end{equation}


边界项可以改写为
\begin{equation}\label{eq:ode-adjoint-boundary-matrix-form}
\begin{aligned}
&\sum_{k,l=1}^{2p}A_{kl}(b)D^{k-1}\phi(b)D^{l-1}\psi(b)\\
&\qquad-\sum_{k,l=1}^{2p}A_{kl}(a)D^{k-1}\phi(a)D^{l-1}\psi(a).
\end{aligned}
\end{equation}
这里当 $k+l>2p+1$ 时 $A_{kl}$ 为零, 且 $A_{k(2p+1-k)}=(-1)^{k-1}a_{2p}$. 由于已假设 $a_{2p}$ 非零, 这意味着矩阵 $\bm{A}$ 非奇异. 现在假设在点 $b$ 处有 $p$ 个线性无关的边界条件
\begin{equation}\label{eq:ode-independent-endpoint-boundary-conditions}
\sum_{j=1}^{2p}b_{ij}D^{j-1}\psi(b)=0,\quad i=1,\ldots,p.
\end{equation}
用 $\bm{u}$ 表示分量为 $u_i=D^{i-1}\psi(b)$ 的 $2p$ 维向量, 用 $\bm{v}$ 表示分量为 $v_i=D^{i-1}\phi(b)$ 的 $2p$ 维向量. 那么 \eqref{eq:ode-independent-endpoint-boundary-conditions} 将 $\bm{u}$ 限制在 $\mathbb{R}^{2p}$ 的一个 $p$ 维子空间 $X$ 内. $X$ 在 $\bm{A}(b)$ 下的像是 $\mathbb{R}^{2p}$ 的一个 $p$ 维子空间 $Y$. 要使 \eqref{eq:ode-adjoint-boundary-matrix-form} 中的第一项对每个满足 \eqref{eq:ode-independent-endpoint-boundary-conditions} 的 $\psi$ 都为零, 充分必要条件是 $\bm{v}$ 位于 $Y$ 的正交补中. 这给出关于 $\phi$ 的一组 $p$ 个边界条件, 称之为伴随边界条件. 在点 $a$ 处也可作类似考察.

作为例子, 设 $L\psi=\psi''''$, 并在每个端点施加边界条件 $\psi+\psi^{\prime}=2\psi+\psi^{\prime\prime}=0$. 此时 $L^*=L$, 而 \eqref{eq:ode-formal-adjoint-boundary-expansion} 化为
\begin{equation}\label{eq:fourth-order-ode-green-identity}
\begin{aligned}
\int_a^b\phi(x)\psi''''(x)\,\mathrm{d}x
&=\int_a^b\phi''''(x)\psi(x)\,\mathrm{d}x\\
&\quad+\phi(b)\psi'''(b)-\phi'(b)\psi''(b)\\
&\quad+\phi''(b)\psi'(b)-\phi'''(b)\psi(b)\\
&\quad-\phi(a)\psi'''(a)+\phi'(a)\psi''(a)\\
&\quad-\phi''(a)\psi'(a)+\phi'''(a)\psi(a).
\end{aligned}
\end{equation}
\eqref{eq:ode-adjoint-boundary-matrix-form} 中的矩阵 $\bm{A}$ 为
\begin{equation}\label{eq:fourth-order-ode-boundary-matrix}
\bm{A}=\begin{pmatrix}0&0&0&1\\0&0&-1&0\\0&1&0&0\\-1&0&0&0\end{pmatrix},
\end{equation}
而向量 $\bm{u}$ 受条件 $u_1+u_2=2u_1+u_3=0$ 约束. 满足这些条件的向量组成的空间 $X$ 的一组基为向量 $(1,-1,-2,0)$ 和 $(0,0,0,1)$. 这些向量在 $\bm{A}$ 下的像为 $(0,2,-1,-1)$ 和 $(1,0,0,0)$. 因此向量 $\bm{v}$ 必须满足条件 $2v_2-v_3-v_4=0$, $v_1=0$, 即伴随边界条件为 $\phi=2\phi^{\prime}-\phi^{\prime\prime}-\phi^{\prime\prime\prime}=0$.
